\section{Conclusion}
\label{sec:conclusion}

An ECG front-end built around an integrated instrumentation amplifier can decide in service, from
26 features it measures on itself, whether it still complies with IEC~60601-2-25, with
\SI{0.9}{\percent} of escapes at \SI{7.8}{\percent} of false rejects; three measurements taking
three seconds come within one point of the balanced accuracy of the complete self-test. Defining
a fault by its effect on the specifications, as alternate test does, instead of by the deviation
of a component reduces false rejects from \SI{37}{\percent} to about \SI{2}{\percent}. The cause
is located with a macro F1 of 0.87 when the answer is an ambiguity group obtained from the
sensitivities of the nominal circuit, and that answer, unlike the component, remains valid for
fault magnitudes absent from training. Non-compliant circuits are separated from electrode
problems; degraded contacts are not, because they fall within the variation between subjects. A
discrete three-op-amp front-end gives a worse picture, so conclusions drawn on benchmark-like
circuits do not carry over to the circuits used in instruments. The next step is validation on
hardware; multiple faults and a per-patient reference for the electrode contact are left for
future work.
